\documentclass[11pt,reqno]{amsart}
\usepackage[T1]{fontenc}
\usepackage[utf8]{inputenc}
\usepackage{lmodern}
\usepackage{amsmath,amssymb,amsthm,mathtools}
\usepackage[colorlinks=true,linkcolor=blue,citecolor=blue]{hyperref}
\usepackage[nameinlink,capitalize]{cleveref}
\newtheorem{theorem}{Theorem}[section]
\newtheorem{lemma}[theorem]{Lemma}
\newtheorem{proposition}[theorem]{Proposition}
\newtheorem{corollary}[theorem]{Corollary}
\theoremstyle{definition}
\newtheorem{definition}[theorem]{Definition}
\newcommand{\E}{\mathbb E}
\newcommand{\R}{\mathbb R}
\newcommand{\F}{\mathcal F}
\newcommand{\norm}[1]{\left\lVert#1\right\rVert}
\newcommand{\abs}[1]{\left\lvert#1\right\rvert}
\title{Physical-Clock Appendix to Paper III:\ Random Roof Time Change and Semi-Markov Theta Recursions}
\author{Qian Qi}
\date{August 29, 2026}
\begin{document}
\maketitle

\section{Collision time and physical time}

Let the predictable moving collision process have branch probabilities
$w_i(A_k)$ and positive branch roofs
\[
 0<\tau_*\le\tau_i(a)\le\tau^*<\infty.
\]
Define
\[
 \bar\tau(a)=\sum_iw_i(a)\tau_i(a).
\]
On collision slow time $s=k\varepsilon^2$, let
\begin{equation}\label{eq:collision-recursion}
 X_{k+1}^\varepsilon=X_k^\varepsilon
 +\varepsilon m_{I_k}(A_k)+\varepsilon^2b(X_k^\varepsilon),
 \qquad A_k=\Theta(X_k^\varepsilon).
\end{equation}
The accumulated physical slow time is
\begin{equation}\label{eq:clock}
 C^\varepsilon(s)=
 \varepsilon^2\sum_{k<s/\varepsilon^2}\tau_{I_k}(A_k),
\end{equation}
with linear interpolation.

\begin{lemma}[Clock martingale decomposition]\label{lem:clock-decomp}
One has
\begin{equation}\label{eq:clock-decomp}
 C^\varepsilon(s)=
 \varepsilon^2\sum_{k<s/\varepsilon^2}\bar\tau(A_k)
 +N^\varepsilon(s),
\end{equation}
where $N^\varepsilon$ is a martingale and
\begin{equation}\label{eq:clock-BDG}
 \E\sup_{s\le S}\abs{N^\varepsilon(s)}^2
 \le C_S\varepsilon^2.
\end{equation}
\end{lemma}

\begin{proof}
Exact predictable innovations give
\[
 \E[\tau_{I_k}(A_k)\mid\F_k]=\bar\tau(A_k).
\]
Thus the centered roof increments form a bounded martingale difference array.
Their scaled quadratic variation on $[0,S]$ is bounded by
\[
 \varepsilon^4(S\varepsilon^{-2}+1)(\tau^*-	au_*)^2
 \le C_S\varepsilon^2.
\]
Doob's inequality gives \eqref{eq:clock-BDG}.
\end{proof}

\begin{theorem}[Uniform physical-clock limit]\label{thm:clock-limit}
Suppose the collision-time slow interpolation converges jointly with its
martingale driver to a continuous process $X_s$.  Then
\begin{equation}\label{eq:clock-limit}
 C^\varepsilon(s)\longrightarrow
 C(s)=\int_0^s\bar\tau(\Theta(X_r))dr
\end{equation}
uniformly on compact intervals in probability.  The limit is strictly
increasing and bi-Lipschitz.  If
\[
 E^\varepsilon(t)=\inf\{s:C^\varepsilon(s)>t\},
 \qquad
 E(t)=C^{-1}(t),
\]
then $E^\varepsilon\to E$ uniformly on compact physical-time intervals in
probability.
\end{theorem}

\begin{proof}
The martingale part vanishes by \cref{lem:clock-decomp}.  Tightness and
continuity of $\bar\tau\circ\Theta$ turn the predictable sum into the Riemann
integral in \eqref{eq:clock-limit}.  The bounds
$\tau_*\le\bar\tau\le\tau^*$ imply
\[
 \tau_*(s-r)\le C(s)-C(r)\le\tau^*(s-r).
\]
Uniform convergence of strictly increasing functions with a common positive
slope bound implies uniform convergence of their right-continuous inverses.
\end{proof}

\section{Time change of the diffusion limit}

Suppose the collision-time limit is
\begin{equation}\label{eq:collision-SDE}
 dX_s=b(X_s)ds+\sigma_c(\Theta(X_s))dW_s,
 \qquad \sigma_c\sigma_c^T=C_{\rm coll}.
\end{equation}
Set $\widetilde X_t=X_{E(t)}$.

\begin{theorem}[Physical-time diffusion]\label{thm:physical-SDE}
The physical-time process is the unique weak solution of
\begin{equation}\label{eq:physical-SDE}
 d\widetilde X_t=
 {b(\widetilde X_t)\over\bar\tau(\Theta(\widetilde X_t))}dt
 +{\sigma_c(\Theta(\widetilde X_t))
   \over\sqrt{\bar\tau(\Theta(\widetilde X_t))}}\,d\widetilde W_t.
\end{equation}
Its covariance is
\begin{equation}\label{eq:physical-covariance}
 \Sigma_{\rm phys}(a)={C_{\rm coll}(a)\over\bar\tau(a)},
\end{equation}
which is exactly the pressure-root covariance of Paper II.  The physical-time
interpolation $X^\varepsilon_{E^\varepsilon(t)}$ converges to
$\widetilde X$ jointly with its time-changed geometric rough lift.
\end{theorem}

\begin{proof}
The inverse clock satisfies
$dE(t)=dt/\bar\tau(\Theta(X_{E(t)}))$.  The continuous martingale
$W_{E(t)}$ has quadratic variation $E(t)$; by the time-change theorem there is
a Brownian motion $\widetilde W$ such that
\[
 dW_{E(t)}={1\over\sqrt{\bar\tau(\Theta(\widetilde X_t))}}d\widetilde W_t.
\]
Substitution in \eqref{eq:collision-SDE} gives
\eqref{eq:physical-SDE}.  The covariance is
\eqref{eq:physical-covariance}.  Joint convergence follows from the uniform
clock and inverse-clock limits and continuity of time change in the uniform
and rough-path topologies for strictly increasing bi-Lipschitz clocks.
\end{proof}

If the collision-time rough limit has a deterministic area anomaly
$\Gamma(a)ds$, physical time changes it to
$\Gamma(a)dt/\bar\tau(a)$.

\section{A branch-dependent physical-time DPP}

Let $h>0$ be the collision scale.  A branch-$i$ collision advances physical
time by $h\tau_i(a)$, slow position by
$hb(x)+\sqrt h\,m_i(a)$, and has probability $w_i(a)$.  Define, for a bounded
continuous space--time function $U$,
\begin{equation}\label{eq:semi-Markov-step}
 (\mathcal S_hU)(t,x)=
 {1\over\vartheta}\log
 \sum_iw_i(\Theta(x))
 \exp\left\{\vartheta
 U\left(t+h\tau_i(\Theta(x)),
 x+hb(x)+\sqrt h\,m_i(\Theta(x))\right)\right\}.
\end{equation}
Starting at $(t,x)$, iterate until the first collision time whose physical clock
reaches or exceeds $T$, and evaluate the terminal payoff at the overshooting
state.  The overshoot is at most $h\tau^*$.

\begin{lemma}[Physical-time consistency]\label{lem:physical-consistency}
For $U\in C_b^{1,3}$,
\begin{align}
 {\mathcal S_hU-U\over h}
 \longrightarrow&\ 
 \bar\tau(\Theta(x))U_t+b\cdot DU
 +{1\over2}C_{\rm coll}(\Theta(x)):D^2U\\
 &+{\vartheta\over2}
 DU^TC_{\rm coll}(\Theta(x))DU
 \label{eq:physical-consistency}
\end{align}
locally uniformly.
\end{lemma}

\begin{proof}
Taylor-expand in both time and space.  The branch mean of the time increment is
$h\bar\tau$.  The centered spatial first-order term vanishes, the spatial
second moment is $hC_{\rm coll}$, and the logarithmic cumulant contributes the
last term.  Mixed time--space and higher terms are $O(h^{3/2})$ because roofs
and branch increments are bounded.
\end{proof}

\begin{theorem}[Physical-time theta equation]\label{thm:physical-HJB}
The semi-Markov terminal recursion generated by
\eqref{eq:semi-Markov-step} converges locally uniformly to the unique bounded
viscosity solution of
\begin{equation}\label{eq:physical-HJB}
 -U_t-{b\over\bar\tau}\cdot DU
 -{1\over2}{C_{\rm coll}\over\bar\tau}:D^2U
 -{\vartheta\over2\bar\tau}
 DU^TC_{\rm coll}DU=0,
 \qquad U(T)=\phi.
\end{equation}
The diffusion coefficient in \eqref{eq:physical-HJB} is precisely
$\Sigma_{\rm phys}=C_{\rm coll}/\bar\tau$.  The same conclusion holds with a
compact control supremum or with the pure lower/upper game operators, provided
the structural hypotheses are uniform in the controls.
\end{theorem}

\begin{proof}
The semi-Markov operator is monotone and cash invariant.  Its physical-time
mesh is nonuniform but lies between $h\tau_*$ and $h\tau^*$, so stability and
the half-relaxed-limit contact argument are unchanged.  Divide the consistency
operator \eqref{eq:physical-consistency} by the positive
$\bar\tau$.  Uniform ellipticity and global coefficient continuity give
comparison.  The terminal overshoot contributes at most the modulus of
continuity of the terminal payoff at space displacement $O(\sqrt h)$ and time
displacement $O(h)$, which tends to zero.  The controlled and game branches
follow by uniform consistency before optimization.
\end{proof}

\section{Compatibility with the full-frequency roof}

The roof values in the physical-clock theorem may be chosen to satisfy the
Diophantine condition of Paper II.  The physical-time homogenization and HJB
use only bounded positivity and smooth parameter dependence; the high-frequency
suspension theorem uses the additional Diophantine separation.  Therefore the
same branchwise roof vector simultaneously carries:

\begin{itemize}
\item the full-frequency moving-family suspension resolvent;
\item the random physical clock and its inverse;
\item the pressure-root covariance;
\item the physical-time rough homogenization;
\item the semi-Markov microscopic theta recursion.
\end{itemize}

\section{Conclusion}

Collision count and physical time are connected by an actual random roof
clock, not by an informal division after the limit.  The clock martingale
vanishes, its predictable part converges, and inverse clocks converge uniformly.
The time-changed diffusion and the branch-dependent microscopic DPP both
produce exactly the covariance $C_{\rm coll}/\bar\tau$ from the pressure-root
theorem, while retaining the same nonlattice roof needed at high frequency.

\bibliographystyle{plain}
\bibliography{references}
\end{document}
